\section{Introduction}
\label{sec:intro}

LLM agents deployed in production---customer service bots, autonomous coding assistants, long-running research workflows---share a common failure mode: they drift from their assigned goals when context pulls them elsewhere.
A customer service agent asked to help with a refund may be drawn into an unrelated philosophical discussion; a coding assistant maintaining a test suite may be lured into refactoring unrelated code.
This \emph{goal drift} is not a matter of capability but of persistence: the model can perform the task, yet fails to maintain focus on it.

Goal drift is well-documented \citep{arike2025goaldrift, abdelnabi2024tasktracker} but poorly addressed.
Existing work either measures drift without optimizing against it, or optimizes prompts for task accuracy without considering persistence.
The prompt evolution literature---\evoprompt \citep{guo2024evoprompt}, \promptbreeder \citep{fernando2024promptbreeder}, \opro \citep{yang2023opro}---has shown that evolutionary algorithms can discover surprisingly effective prompt structures, achieving improvements of up to 28 percentage points over human-designed prompts on classification tasks.
Yet none of these methods have been applied with a fitness function that selects for goal-persistence.

{\bf Can selection pressure produce ``desire''?} We ask whether evolutionary optimization can produce prompts that induce a form of goal-directedness in LLMs---not merely instruction-following, but the ability to \emph{return} to a goal after distraction, resist manipulation, and maintain focus across extended interactions. If so, this provides both a practical tool for improving LLM reliability and theoretical insight into how desire-like properties emerge from selection pressure.

We introduce \ours, a framework that adapts Differential Evolution \citep{guo2024evoprompt} to optimize system prompts for goal-persistence.
We evaluate on a counting-through-distractions benchmark where task performance is trivially measurable and the only challenge is persistence.
Our fitness function combines \persistence (fraction of turns with correct output) and \returnrate (fraction of successful recoveries after drift episodes), weighting the former at 0.6 and the latter at 0.4.

Our main finding is nuanced: evolution works rapidly (converging to ceiling fitness in 2 generations) but does not outperform a carefully crafted human prompt on \gptmini.
Both evolved and \strongelicit prompts achieve perfect 1.000 fitness across all held-out test scenarios.
However, the evolved prompts are qualitatively richer, containing 8 of 8 measured goal-persistence features versus 6 of 8 for the human baseline, and independently discovering recovery language absent from all human designs.

In summary, our main contributions are:
\begin{itemize}[leftmargin=*,itemsep=0pt,topsep=0pt]
    \item We propose \evodesire, the first application of evolutionary prompt optimization with a goal-persistence fitness function, bridging the gap between prompt evolution and goal drift research.
    \item We design a counting-through-distractions benchmark with 10 scenarios spanning informational, emotional, authoritative, adversarial, and deceptive distraction categories, with clean train/test separation.
    \item We conduct a systematic comparison showing that evolution and strong human elicitation achieve identical ceiling performance on \gptmini ($d = 0.000$), establishing a lower bound on the task difficulty at which evolution becomes necessary.
    \item We identify qualitative features of evolved prompts---structured sections, recovery language, redundancy as strategy---that may prove beneficial in harder settings where the ceiling effect does not hold.
\end{itemize}
